%% 29_arch.tex -- Chapter 29: The Full Architecture and Its Budgets
%% Status: OUTLINE (build stage 6). Replace the outline with the full chapter when it is written.
\chapter{The Full Architecture and Its Budgets}\label{ch:arch}
\chapterlead{Every chapter so far has answered one budget question. This chapter puts the answers together on one chip.}

\begin{specnote}[Chapter status: outline]
This chapter is an outline. It will be written in build stage 6 (Part V, Lab 14, the capstone and the appendices). The topics, the design question and the labs below are fixed; the derivations, worked examples and figures are still to come.
\end{specnote}

\begin{objectives}
\begin{itemize}
  \item assemble the noise, CMRR, power, area and data-rate budgets of \OB{};
  \item choose SAR or $\Delta\Sigma$ per mode;
  \item add the optional bioimpedance channel;
  \item draw the top-level floorplan.
\end{itemize}
\end{objectives}

\begin{insight}[Design question]
Does the assembled \OB{} chip meet every number of its specification (\cref{tab:ch1-spec}), and where is the margin?
\end{insight}

\section{Noise budget}\label{sec:ch29-1}
From electrode noise (\cref{ch:electrodes}) through the LNA, PGA and ADC to the \qty{1.5}{\uVrms} EEG target.

\section{CMRR budget}\label{sec:ch29-2}
Electrode imbalance, amplifier mismatch and right-leg drive against the \qty{90}{\deci\bel} target.

\section{Power, area and data rate}\label{sec:ch29-3}
The \qty{10}{\micro\watt} per channel front-end budget, converter energy, and the output data rate per mode.

\section{SAR or \texorpdfstring{$\Delta\Sigma$}{Delta-Sigma} per mode}\label{sec:ch29-4}
The routing of \cref{fig:ch1-chain} justified by the budgets.

\section{The bioimpedance channel}\label{sec:ch29-5}
The optional respiration channel of \cref{ch:bioz}.

\section{Floorplan}\label{sec:ch29-6}
Placing eight channels, two converters, references and the digital back end.

\paragraph{Read with this chapter} Verma et al.~\cite{verma2010}, Yazicioglu, Van Hoof and Puers~\cite{yazicioglubook}.
